% ============================================================
% analy 报告 — QUDA multigrid 算法结构与思路分析
% 编译: xelatex 两遍; 交付前 Overfull 与 Float too large 计数必须为 0
% ============================================================
\documentclass[UTF8]{ctexart}
\usepackage[margin=2.5cm]{geometry}
\usepackage{graphicx}
\usepackage{amsmath}
\usepackage{microtype}
\usepackage{listings}
\usepackage{xcolor}
\usepackage{booktabs}
\usepackage{longtable}
\usepackage{xltabular}
\usepackage{tabularx}
\usepackage{hyperref}
\usepackage{fancyhdr}
\usepackage[most]{tcolorbox}
\usepackage{titlesec}

% ===== 色板 =====
\definecolor{accent}{HTML}{1F4E79}
\definecolor{evidence}{HTML}{1E8449}
\definecolor{reference}{HTML}{8C6D1F}
\definecolor{conclusion}{HTML}{8B1A1A}
\definecolor{codebg}{HTML}{F4F6F7}
\definecolor{struct}{HTML}{3B3B98}
\definecolor{relation}{HTML}{0E7C7B}
\definecolor{idea}{HTML}{B03A2E}
\definecolor{warn}{HTML}{D35400}
\definecolor{hl}{HTML}{FDF6E3}

\setlength{\emergencystretch}{3em}
\hypersetup{colorlinks=true,linkcolor=accent,urlcolor=relation}

\titleformat{\section}{\Large\bfseries\color{accent}}{\thesection}{1em}{}
\titleformat{\subsection}{\large\bfseries\color{struct}}{\thesubsection}{1em}{}
\titleformat{\subsubsection}{\normalsize\bfseries\color{struct!70!black}}{\thesubsubsection}{1em}{}

\newtcolorbox{conclbox}{breakable,colback=conclusion!6,colframe=conclusion,title=结论}
\newtcolorbox{refbox}{breakable,colback=reference!6,colframe=reference,title=参考背景}
\newtcolorbox{ideabox}{breakable,colback=idea!6,colframe=idea,title=项目思路}
\newtcolorbox{warnbox}{breakable,colback=warn!6,colframe=warn,title=局限与风险}
\newtcolorbox{keybox}{breakable,colback=hl,colframe=accent,title=要点}

\lstset{
  basicstyle=\ttfamily\footnotesize,
  backgroundcolor=\color{codebg},
  frame=single,
  language=C++,
  firstnumber=auto,
  keywordstyle=\color{accent}\bfseries,
  commentstyle=\color{evidence!70!black},
  stringstyle=\color{struct},
  breaklines=true,
  breakatwhitespace=false,
  postbreak=\mbox{\textcolor{accent}{$\hookrightarrow$}\space},
  keepspaces=true,
  columns=flexible,
  numbers=left,numberstyle=\tiny\color{struct!60!black},
}

\title{QUDA 1.1.0 自适应多重网格求解器\\整体算法思路与结构分析}
\author{opencode analy}
\date{\today}

\begin{document}
\maketitle

\begin{abstract}
本报告以只读方式分析 QUDA 1.1.0（格点 QCD GPU 库）中自适应多重网格（Adaptive Multigrid, MG）求解器的整体算法思路与代码结构。分析范围覆盖仓库全部 806 个非 git 文件（精读 5 个、浏览 20 余个、其余按目录登记仅索引），主题聚焦于 \texttt{\nolinkurl{include/multigrid.h}}、\texttt{\nolinkurl{lib/multigrid.cpp}}、\texttt{\nolinkurl{include/transfer.h}}、\texttt{\nolinkurl{lib/transfer.cpp}} 及粗算子构建内核族。核心结论：(1) QUDA MG 是\textbf{聚合型代数多重网格}——粗空间由数值生成的近零模向量张成，而非几何光滑函数；(2) 其骨架是\textbf{每层一个递归 MG 实例}的组合结构：Transfer（$R=P^\dagger$ 正交投影）+ Galerkin 粗算子 $D_c=RDP$（粗链接 $Y$、粗 clover $X$）+ 前/后平滑器 + 粗网格求解器；(3) ``自适应''体现为 setup 阶段的迭代自举——零空间向量经平滑求解生成后逐级限制下传并重建算子；(4) 数值正确性由内建 \texttt{verify()} 五项投影恒等式校验保障。报告按``结构—关系—思路''三视角组织，主题分析节采用 A（代码工作）15 步框架展开。
\end{abstract}

\tableofcontents

% ================= 1 仓库概览 =================
\section{仓库概览}

\subsection{目录结构与规模}

QUDA 1.1.0 为纯 C++/CUDA 库仓库，当前分支 main（HEAD \texttt{1da4469}）。顶层目录实测如下表。

\begin{xltabular}{\textwidth}{llX}
\toprule
\textcolor{struct}{目录} & 文件数 & \textcolor{struct}{职责} \\
\midrule
\endhead
\texttt{\nolinkurl{lib/}} & 326 & 核心实现层：Dirac 算子、求解器、multigrid、blas、通信 \\
\texttt{\nolinkurl{include/}} & 310 & 接口与模板头：公共 API、内核参数结构、模板声明 \\
\texttt{\nolinkurl{tests/}} & 144 & 测试入口：invert\_test、multigrid\_benchmark/evolve、参数工具 \\
\texttt{\nolinkurl{docs/}} & 6 & 本库历史分析报告（参考知识库，见 \S9 类型列``参考''） \\
\texttt{\nolinkurl{doc/}}、\texttt{\nolinkurl{cmake/}}、\texttt{\nolinkurl{ci/}}、\texttt{\nolinkurl{jekins/}} & 11 & Doxygen 配置与构建/CI 基础设施 \\
根目录散文件 & 5 & README.md、NEWS、CMakeLists.txt、LICENSE 等 \\
\bottomrule
\caption{顶层目录职责与规模（数据来源：find 实测）}
\end{xltabular}

\begin{table}[htbp]\centering
\begin{tabular}{ll}
\toprule
指标 & 实测值 \\
\midrule
C++ 源文件（.cpp/.cu/.cuh/.h） & 183 / 154 / 142 / 213 \\
脚本与文档（.sh/.py/.md） & 7 / 2 / 4 \\
git 标签数（test 系列） & $\ge$20（test15\_1$\sim$test20） \\
MG 递归深度上限 & 5 层（\texttt{\nolinkurl{include/quda_constants.h:50}}） \\
\bottomrule
\end{tabular}
\caption{仓库实测统计（数据来源：find/wc/git 实测）}
\end{table}

README 明确定位该组件：``QUDA includes an implementations of adaptive multigrid for the Wilson, clover-improved, twisted-mass and twisted-clover fermion actions''\textcolor{evidence}{\texttt{\nolinkurl{README.md:28-30}}}。NEWS 记录其演进：``Initial version of adaptive multigrid fully implemented into QUDA''\textcolor{evidence}{\texttt{\nolinkurl{NEWS:207}}}，至 1.1.0 已支持 tensor core 加速的 setup\textcolor{evidence}{\texttt{\nolinkurl{NEWS:33-35}}}与 staggered 多重网格\textcolor{evidence}{\texttt{\nolinkurl{NEWS:29-30}}}。

\subsection{覆盖登记表（饱和声明）}

四条遍历路径交叉走遍全库，覆盖表闭合如下（按目录汇总计数）：

\begin{xltabular}{\textwidth}{lXX}
\toprule
路径 & 三态分配 & 证据命中 \\
\midrule
\endhead
lib/(305 文件) & 精读 5：multigrid.cpp、transfer.cpp、dirac\_coarse.cpp、coarse\_op.cuh、solver.cpp(MG 段)；浏览 10 余个 mg 文件族；其余仅索引（blas/comm/HMC 等非主题文件，历史报告已覆盖其结构） & \textcolor{evidence}{\texttt{\nolinkurl{multigrid.cpp:1131}}} 等 \\
include/(310) & 精读 2：multigrid.h、transfer.h；浏览 8 个头（quda.h 参数段、enum\_quda.h、multigrid\_helper.cuh 与 kernels/ 下 prolongator、restrictor、coarse\_op\_kernel、block orthogonalize 四内核头）；其余仅索引 & \textcolor{evidence}{\texttt{\nolinkurl{multigrid.h:263}}} 等 \\
tests/(144) & 浏览 4：invert\_test.cpp(MG 段)、utils/set\_params.cpp 的 setMultigridParam、multigrid\_benchmark 与 multigrid\_evolve 两测试头部；其余仅索引 & set\_params.cpp:361 等 \\
docs/(6)、根散文件(5)、doc/cmake/ci/jenkins(11) & docs/ 三份 PDF 经 pdftotext 检索命中 multigrid 主题；README.md、NEWS 浏览；其余仅索引 & README.md:28 等 \\
\bottomrule
\caption{覆盖登记表闭合汇总（数据来源：find/grep/pdftotext 实测）}
\end{xltabular}

% ================= 2 主体结构 =================
\section{主体结构}

multigrid 子系统在仓库中呈六层结构：API 入口层把用户参数翻译为内部对象树；算法驱动层 \texttt{MG} 类是唯一持有递归结构的实体；层间转移层与粗算子层分别实现``怎么上下''与``下面是什么''；求解器工厂层提供平滑器与粗解器；测试层给出典型调用序。

\begin{xltabular}{\textwidth}{lp{3.0cm}p{3.6cm}X}
\toprule
\textcolor{struct}{层次} & 成员文件 & 关键符号 & \textcolor{struct}{职责一句话} \\
\midrule
\endhead
入口层 & \texttt{\nolinkurl{lib/interface_quda.cpp}} & \texttt{\nolinkurl{newMultigridQuda}}、\texttt{\nolinkurl{multigrid_solver}} & 校验参数、创建细格 Dirac 三件套与零空间场 B，组装 MGParam 并启动 MG 树 \\
接口声明 & \texttt{\nolinkurl{include/quda.h}} & \texttt{\nolinkurl{QudaMultigridParam}} & 逐层参数表（block、n\_vec、setup、smoother、coarse\_solver、cycle\_type 等 40 余项） \\
算法驱动 & \texttt{\nolinkurl{lib/multigrid.cpp}}、\texttt{\nolinkurl{include/multigrid.h}} & \texttt{class MG : public Solver} & setup 递归（reset）与求解递归（operator() 即 V-cycle） \\
层间转移 & \texttt{\nolinkurl{lib/transfer.cpp}} 与 kernels/prolongator、restrictor、block orthogonalize 内核头 & \texttt{\nolinkurl{Transfer::P}} / \texttt{\nolinkurl{Transfer::R}} & 聚合映射、块正交化生成 V 场、实现 $P$ 与 $R=P^\dagger$ \\
粗算子 & \texttt{\nolinkurl{lib/dirac_coarse.cpp}}、lib/coarse\_op*.cuh、kernels/coarse op kernel.cuh & \texttt{\nolinkurl{CoarseOp}}、\texttt{\nolinkurl{CoarseCoarseOp}}、\texttt{\nolinkurl{calculateYhat}} & 构造粗链接 $Y$、粗 clover $X$ 及预条件形式 $\hat Y=X^{-1}Y$ \\
求解器工厂 & \texttt{\nolinkurl{lib/solver.cpp}}、inv\_mr、inv\_gcr 等求解器文件 & \texttt{\nolinkurl{Solver::create}} & 平滑器（MR/GCR + Schwarz）与粗解器实例化 \\
测试层 & \texttt{\nolinkurl{tests/invert_test.cpp}}、\texttt{\nolinkurl{tests/utils/set_params.cpp}} & \texttt{\nolinkurl{setMultigridParam}} & 典型配置与调用生命周期演示 \\
\bottomrule
\caption{multigrid 子系统分层（数据来源：全库索引实测）}
\end{xltabular}

类关系上（\textcolor{evidence}{\texttt{\nolinkurl{include/multigrid.h:263}}}）：\texttt{MG} 继承抽象基类 \texttt{Solver}，其成员 \texttt{coarse} 指向下一层 \texttt{MG} 实例（\textcolor{evidence}{\texttt{\nolinkurl{multigrid.h:287}}}），\texttt{transfer} 指向 \texttt{Transfer}（\textcolor{evidence}{\texttt{\nolinkurl{multigrid.h:270}}}），而 \texttt{presmoother}、\texttt{postsmoother}、\texttt{coarse\_solver} 均为 \texttt{Solver} 工厂产物（\textcolor{evidence}{\texttt{\nolinkurl{multigrid.h:276--290}}}）——即\textbf{MG 树本身就是一棵``组合模式''的对象树}，递归深度由 \texttt{param.level}$<$\texttt{Nlevel} 控制（\textcolor{evidence}{\texttt{\nolinkurl{multigrid.cpp:95}}}）。

% ================= 3 各部分关系 =================
\section{各部分关系}

各部件通过``构造期依赖注入 + 运行期递归调用''连接。以下文本链中箭头表示``创建/调用''关系，均附实测参考源。

\begin{keybox}
\textbf{setup 依赖主链（构造期）：}\\
\textcolor{relation}{\texttt{newMultigridQuda} → \texttt{multigrid\_solver} → \texttt{MG(level 0)::reset()} → \texttt{Transfer}(B) → \texttt{BlockOrthogonalize} → $V$ 场}\\
\textcolor{relation}{\hspace*{4em}→ \texttt{createCoarseDirac()} → \texttt{CoarseOp}/\texttt{CoarseCoarseOp} → $(Y,X)$ → \texttt{calculateYhat} → $(\hat Y,X^{-1})$}\\
\textcolor{relation}{\hspace*{4em}→ \texttt{createSmoother()} → \texttt{new MG(level+1)} 递归 → \texttt{createCoarseSolver()} → \texttt{PreconditionedSolver(GCR)}}\\[0.5em]
\textbf{solve 依赖主链（运行期）：}\\
\textcolor{relation}{外层 GCR(\texttt{invertQUDA}) → \texttt{MG::operator()} → pre-smooth → \texttt{Transfer::R} → \texttt{coarse\_solver} → \texttt{Transfer::P} → post-smooth}
\end{keybox}

三条关键引用关系：（1）\texttt{DiracCoarse::createCoarseOp} 只接受聚合粗化，``Coarse operators only support aggregation coarsening''\textcolor{evidence}{\texttt{\nolinkurl{lib/dirac_coarse.cpp:464-465}}}；（2）外层求解器工厂在 \texttt{inv\_type==GCR} 且带预条件子时取出 MG 实例包装为 \texttt{GCR(mat, *(mg), ...)}\textcolor{evidence}{\texttt{\nolinkurl{lib/solver.cpp:69-76}}}；（3）粗算子构建需要细算子提供 \texttt{kappa/mass/mu} 与 \texttt{mu\_factor} 差值（\textcolor{evidence}{\texttt{\nolinkurl{multigrid.cpp:371--376}}}），体现逐层 twisted-mass 结构的一致传递。演进关系上，NEWS 显示 staggered KD 粗化类型与 MMA 加速为后期增量（\texttt{\nolinkurl{NEWS:29-35}}），而核心 AGGREGATE 路径自初版即存在（\texttt{\nolinkurl{NEWS:207}}）。

% ================= 4 项目思路 =================
\section{项目思路}

\subsection{为什么是``聚合代数''多重网格}

Wilson-Dirac 方程 $D\psi=b$ 在接近临界质量时条件数恶化，Krylov 方法收敛由最小特征模（长程物理，如 pion 极点）主导。经典几何多重网格假定解光滑，但规范场构型随机粗糙、且费米子矩阵含 $\gamma_\mu$ 反对易结构，光滑性假设不成立。QUDA 的选择是：让\textbf{数值生成的近零模向量}自己定义粗空间——把每个 $2^4$ 几何块（aggregate）内的格点捆绑为一个粗格点，块内以近零模向量为基做正交投影。这正是代数多重网格（AMG）思想，而``自适应''指近零模不是先验给定的，而是 setup 阶段迭代求出的。

\begin{ideabox}
整个子系统回答一个问题：\textbf{如何在 GPU 上把病态的格点 Dirac 反演拆解为``高频局部平滑 + 低频粗网格精确解''}。为此它把``误差分频''物理直觉落实为三个工程构件——正交投影式转移算子（Transfer）、Galerkin 投影粗算子（$D_c=RDP$，保 Hermiticity 与谱结构）、以及逐级递归的 MG 对象树；并把``近零空间从哪来''做成可插拔的自适应 setup 流程。
\end{ideabox}

\subsection{设计意图的四个落点}

(1)\textbf{可验证性内建}：\texttt{MG::verify()} 把五项投影恒等式写成运行时检查（\textcolor{evidence}{\texttt{\nolinkurl{multigrid.h:433--443}}} 文档注释逐项列出）；(2)\textbf{GPU 优先}：混合精度链 double→single→half→quarter（\textcolor{evidence}{\texttt{\nolinkurl{multigrid.h:43--69}}} 的 \texttt{instantiatePrecisionMG}）、MMA tensor core 开关贯穿 setup/dslash/transfer 三处（\textcolor{evidence}{\texttt{\nolinkurl{multigrid.h:169--176}}}）；(3)\textbf{血统可溯}：setup 类型 \texttt{QUDA\_TEST\_VECTOR\_SETUP} 分支注释直书``DDalphaAMG test vector idea solving against the vector''\textcolor{evidence}{\texttt{\nolinkurl{lib/multigrid.cpp:1386-1388}}}，表明其方法学源自 DDalphaAMG 项目；(4)\textbf{复用友好}：\texttt{multigrid\_solver} 以不透明指针形态存活于多次 solve 之间（\textcolor{evidence}{\texttt{\nolinkurl{multigrid.h:712--717}}}），HMC 分子动力学演化中经 \texttt{updateMultigridQuda} 仅重置算子不重建零空间（\textcolor{evidence}{\texttt{\nolinkurl{multigrid.h:387}}} \texttt{reset(bool refresh)}）。

\subsection{典型工作流}

一次典型反演走遍全库：测试程序填 \texttt{QudaMultigridParam}（默认 RECURSIVE cycle + AGGREGATE 转移，\textcolor{evidence}{\texttt{\nolinkurl{tests/utils/set_params.cpp:471-474}}}）→ \texttt{newMultigridQuda} 触发 setup 递归（生成/加载零空间、逐级建 $V$、$(Y,X)$、平滑器、粗解器）→ 主反演 \texttt{invertQUDA} 外层 GCR 每步调 \texttt{MG::operator()} 做 V-cycle → 结束 \texttt{destroyMultigridQuda} 递归析构（tests/invert\_test.cpp:239、385）。

% ================= 5 主题分析 =================
\section{主题分析：multigrid 整体算法思路与结构}

被分析对象是一项具体的数值求解工作（代码+物理交叉），按规则选 A 框架为主体，物理假设并入相应步骤。判定依据：主题词``multigrid 算法……整体算法思路与结构''指向软件系统的完整工作流重现。

\subsection{A1 目的与前期准备}

目的：为 Wilson/clover/twisted(-clover)/staggered 族 Dirac 算子 $D\psi=b$ 提供随规模弱标度的预条件子。前置条件：已初始化 QUDA、规范场（及 clover 场）已上传 GPU；用户须逐层给定聚合尺寸 \texttt{geo\_block\_size}、粗自由度 \texttt{n\_vec}、平滑器类型与迭代数等参数（\texttt{\nolinkurl{include/quda.h:638-830}}）。层数上限硬编码为 \texttt{QUDA\_MAX\_MG\_LEVEL=5}\textcolor{evidence}{\texttt{\nolinkurl{include/quda_constants.h:46-50}}}。

\subsection{A2 输入是什么}

输入分三类。（1）\textbf{物理场输入}：规范链接 $U_\mu(x)$ 与可选 clover 场，由细格 Dirac 算子持有（interface\_quda.cpp:2860）。（2）\textbf{参数输入}：\texttt{QudaMultigridParam} 逐层数组——几何块尺寸、自旋块尺寸、零空间维数 \texttt{n\_vec}、块正交次数 \texttt{n\_block\_ortho}、setup 求解器 \texttt{setup\_inv\_type} 与迭代预算 \texttt{num\_setup\_iter/setup\_maxiter}、平滑器 \texttt{smoother} 与 \texttt{nu\_pre/nu\_post}、粗解器三元组 \texttt{coarse\_solver[\_tol,\_maxiter]}、循环类型 \texttt{cycle\_type}（\textcolor{evidence}{\texttt{\nolinkurl{include/quda.h:636-830}}}）。（3）\textbf{初始粗空间}：零空间向量组 $B$ 的来源开关——随机初始化后数值生成（默认）、文件加载 \texttt{vec\_load}、特征求解器 \texttt{use\_eig\_solver}、或自由场解析构造 \texttt{buildFreeVectors}（\textcolor{evidence}{\texttt{\nolinkurl{multigrid.cpp:38--62}}} 四分支）。测试端默认值：\texttt{n\_vec} 缺省 24、\texttt{n\_level=mg\_levels}、\texttt{geo\_block\_size} 缺省 4（\textcolor{evidence}{\texttt{\nolinkurl{tests/utils/set_params.cpp:435-460}}}）。

\subsection{A3 输出应该是什么}

输出有三层：（1）解向量 $x$ 满足残差降至 \texttt{tol}；（2）一个可在多次反演间存活的预条件子实例 \texttt{void*}（\texttt{newMultigridQuda} 返回值，interface\_quda.cpp:2929--2939）；（3）副产品——各级零空间可持久化（\texttt{saveVectors}/\texttt{loadVectors}，\textcolor{evidence}{\texttt{\nolinkurl{multigrid.cpp:1226--1263}}}），粗格 deflation 空间可跨 solver 保留（\textcolor{evidence}{\texttt{\nolinkurl{multigrid.cpp:545--548}}}、674--684）。正确性判据：\texttt{verify()} 五项恒等式偏差低于精度相关容差（单精度 $2\times10^{-3}$、半/四分之一精度 $5\times10^{-2}$，\textcolor{evidence}{\texttt{\nolinkurl{multigrid.cpp:754--760}}}）。

\subsection{A4 整体框架：递归对象树与两阶段生命周期}

MG 的结构可用一句话概括：\textbf{每层一个 \texttt{MG} 实例，实例内聚齐``上下来去''的全部构件}。构造函数只做两件事：初始化零空间来源（\textcolor{evidence}{\texttt{\nolinkurl{multigrid.cpp:36--64}}}）然后调 \texttt{reset()}；\texttt{reset()} 才是 setup 的真正主体——它自顶向下递归（\textcolor{evidence}{\texttt{\nolinkurl{multigrid.cpp:95--186}}}）：建 Transfer、建粗 Dirac、建平滑器，再以 \texttt{new MG(*param\_coarse)} 进入下一层（\textcolor{evidence}{\texttt{\nolinkurl{multigrid.cpp:172--178}}}），回溯时再包一层 \texttt{PreconditionedSolver} 作为粗解器（\textcolor{evidence}{\texttt{\nolinkurl{multigrid.cpp:653--667}}}）。

\begin{xltabular}{\textwidth}{llX}
\toprule
\textcolor{struct}{阶段} & 入口方法 & \textcolor{struct}{动作序列} \\
\midrule
\endhead
Setup（一次性） & \texttt{MG::MG} → \texttt{reset()} & 零空间来源选择 → Transfer 创建/复位 → 粗场 $(r\_coarse,x\_coarse,B\_coarse)$ 分配 → \texttt{createCoarseDirac} → \texttt{createSmoother} → 递归下一层 → \texttt{createCoarseSolver} → \texttt{verify()}（\textcolor{evidence}{\texttt{\nolinkurl{multigrid.cpp:72--197}}}） \\
Solve（每步） & \texttt{MG::operator()} & pre-smooth $\nu_{pre}$ 次 → 残差限制 $R$ → 粗解 → 修正延拓 $P$ → post-smooth $\nu_{post}$ 次（\textcolor{evidence}{\texttt{\nolinkurl{multigrid.cpp:1131--1224}}}） \\
Reset（演化复用） & \texttt{reset(true)} & 刷新零空间（\texttt{setup\_maxiter\_refresh}）后重建算子，跳过向量重生成（\textcolor{evidence}{\texttt{\nolinkurl{multigrid.cpp:89--92}}}） \\
\bottomrule
\caption{两阶段生命周期总览（数据来源：源码实测）}
\end{xltabular}

\subsection{A5 关键细节一：转移算子的数学——$R=P^\dagger$}

聚合粗化的核心约定：\textbf{限制就是延拓的厄米共轭}。延拓分两步（\textcolor{evidence}{\texttt{\nolinkurl{prolongator.cuh:141--150}}}）：先把粗向量按自旋映射复制到细自由度（\texttt{prolongate}），再用 $V$ 场做色空间旋转（\texttt{rotateFineColor}）：
\begin{equation}
\eta^f_i(x,s) \;=\; \sum_{j=0}^{N_{vec}-1} V_{ij}(x,s)\,\eta^c_j(x_c,\,s_c),
\label{eq:prolongate}
\end{equation}
其中 $V_{ij}(x,s)=\langle \phi_j(x,s)\,,\,e_i\rangle$ 是块正交化后的近零模分量，$s_c=s/\text{spin\_bs}$ 为手征块索引。\eqref{eq:prolongate} 的实现正是 \texttt{rotateFineColor} 中的 \texttt{cmac(arg.v(...), in[...], partial)} 循环\textcolor{evidence}{\texttt{\nolinkurl{include/kernels/prolongator.cuh:106-116}}}。反向的限制核则做块内归约投影：
\begin{equation}
\eta^c_j(x_c,s_c) \;=\; \sum_{x\in\text{agg}(x_c)}\sum_{s,c}\;\overline{V_{cj}(x,s)}\;\eta^f_c(x,s)\;\equiv\;(P^\dagger\eta^f)_j,
\label{eq:restrict}
\end{equation}
对应 \texttt{rotateCoarseColor} 中 \texttt{cmac(conj(arg.v(...)), in(...), ...)} 与块规约\textcolor{evidence}{\texttt{\nolinkurl{include/kernels/restrictor.cuh:107-116}}}。因为 $V$ 在每个块内正交归一（\texttt{BlockOrthogonalize}，Gram-Schmidt 重复 \texttt{n\_block\_ortho} 次，\textcolor{evidence}{\texttt{\nolinkurl{transfer.cpp:152--164}}}），有精确恒等式 $RP=P^\dagger P=I_c$——这就是 verify 第二项检查的内容。

$V$ 场的存储是个巧妙的工程压缩：把 $N_{vec}$ 个近零模打包成 \texttt{nColor=Ncolor*Nvec} 的``矩阵场''（\textcolor{evidence}{\texttt{\nolinkurl{transfer.cpp:128}}}），使单次访存即可取得某点的全部基分量；粗格点坐标经 \texttt{fine\_to\_coarse} 查找表获得，该表由 \texttt{createGeoMap} 逐点整除块尺寸生成并以排序建立逆表 \texttt{coarse\_to\_fine}\textcolor{evidence}{\texttt{\nolinkurl{lib/transfer.cpp:215-237}}}。聚合尺寸约束 $N_{vec}\le$ 块内自由度总数在构造期强制（\textcolor{evidence}{\texttt{\nolinkurl{transfer.cpp:77--84}}}）。

\subsection{A5 关键细节二：Galerkin 粗算子 $D_c=RDP$ 与规范协变性}

粗算子不是任意近似，而是严格的 Galerkin 投影 $D_c=P^\dagger D P$。构建分计算步（\texttt{COMPUTE\_*} 枚举，coarse\_op.cuh:10--27）：UV/LV（链接乘 $V$）→ AV（clover 乘 $V$）→ VUV（块稀疏收缩 $\sum_x (\overline{AV})\,U\,V$，其中 $\gamma_\mu$ 投影的正负号决定前后向链接合并，coarse\_op\_kernel.cuh:1074--1114）→ COARSE\_CLOVER → DIAGONAL/REVERSE\_Y 收尾。产出两个粗场：粗链接 $Y_\mu(x_c)$ 与粗 clover $X(x_c)$，随后 \texttt{calculateYhat} 生成预条件形式 $\hat Y=X^{-1}Y$ 与 $X^{-1}$（\textcolor{evidence}{\texttt{\nolinkurl{multigrid.h:689--710}}}）。

\textcolor{idea}{格点 MG 区别于普通 PDE MG 的关键是\textbf{规范协变}：近零模在不同格点携带不同规范相位，若直接平均会破坏规范结构。QUDA 的做法是在收缩前把 $V$ 用平行输运搬到相邻聚合（UV 步中的 $U_\mu(x)V(x+\mu)$），使 $Y$ 成为真正的粗格``规范场''，从而粗算子自动继承规范协变性与 $\gamma_5$ 厄米结构。}verify 第三项用随机向量数值检验 $0=(D_c-P^\dagger DP)$\textcolor{evidence}{\texttt{\nolinkurl{lib/multigrid.cpp:883-999}}}，twisted-mass 的逐层 $\mu$ 缩放误差还能被解析扣除（\textcolor{evidence}{\texttt{\nolinkurl{multigrid.cpp:980--991}}}）。

\subsection{A5 关键细节三：自适应零空间生成（``自适应''之所在)}

\texttt{generateNullVectors}（\textcolor{evidence}{\texttt{\nolinkurl{multigrid.cpp:1275--1473}}}）是自适应机制的引擎。流程：（1）以当前（可能尚无粗层的）平滑算子构造 setup 求解器——默认 \texttt{setup\_inv\_type=GCR} 且以自身 MG 为预条件子做一次加性 Schwarz 平滑（\textcolor{evidence}{\texttt{\nolinkurl{multigrid.cpp:1340--1358}}}），即\textbf{自举}：粗层尚未建好时先用平滑器顶替；（2）对 $N_{vec}$ 个随机初始向量批量求解 $M^\dagger M x=0$ 型近零问题（实际是以 $b=0$、随机初值的平滑下降，或 TEST\_VECTOR\_SETUP 模式的 DDalphaAMG 式测试向量法，\textcolor{evidence}{\texttt{\nolinkurl{multigrid.cpp:1385--1413}}}）；（3）pre/post 全局正交归一（\textcolor{evidence}{\texttt{\nolinkurl{multigrid.cpp:1370--1381}}}、1416--1427）；（4）外层重复 \texttt{num\_setup\_iter} 轮，每轮结束后置 \texttt{resetTransfer=true} 并递归 \texttt{reset()}——用改进后的 $B$ 重建本级 $V$、限制下传 $B_{coarse}$、并命令粗层同样重建（\textcolor{evidence}{\texttt{\nolinkurl{multigrid.cpp:1429--1451}}}）。这一``生成→投影→重建''闭环就是文献所称 adaptive setup：粗空间逐步逼近真实近零模张成的低频子空间。

\subsection{A5 关键细节四：V-cycle 主循环}

\texttt{MG::operator()}（\textcolor{evidence}{\texttt{\nolinkurl{multigrid.cpp:1131--1224}}}）是误差分频思想的直接实现。细格残差先交给 pre-smoother（其 \texttt{return\_residual=true}，直接返回平滑后残差避免重复计算，\textcolor{evidence}{\texttt{\nolinkurl{multigrid.cpp:282}}}），残差经 \texttt{transfer->R} 下行，粗解器递归求解粗修正，结果经 \texttt{transfer->P} 上行并与细解累加（\textcolor{evidence}{\texttt{\nolinkurl{multigrid.cpp:1201--1206}}}），最后 post-smoother 清理延拓引入的高频污染。奇偶预条件的解类型一致性（\texttt{smoother\_solver\_uniform}，\textcolor{evidence}{\texttt{\nolinkurl{multigrid.cpp:1148--1158}}}）决定是否需要显式重算残差。最底层（\texttt{level==Nlevel-1}）退化为纯平滑解（\textcolor{evidence}{\texttt{\nolinkurl{multigrid.cpp:1215--1221}}}），可选配粗格 deflation（\textcolor{evidence}{\texttt{\nolinkurl{multigrid.cpp:584--623}}}）。

循环类型枚举为 \texttt{\{VCYCLE, FCYCLE, WCYCLE, RECURSIVE\}}\textcolor{evidence}{\texttt{\nolinkurl{include/enum_quda.h:178-184}}}；但 \texttt{createCoarseSolver} 只区分两条路径——VCYCLE 且非倒数第二层时粗解器直接别名下层 MG 实例（\textcolor{evidence}{\texttt{\nolinkurl{multigrid.cpp:569--572}}}），RECURSIVE 或底层则包 \texttt{PreconditionedSolver}（\textcolor{evidence}{\texttt{\nolinkurl{multigrid.cpp:573--667}}}）。\textcolor{relation}{推断}：FCYCLE/WCYCLE 由外层 PreconditionedSolver 的 \texttt{precondition\_cycle} 参数以多轮 V-cycle 复合实现，代码未见独立 F/W 分支逻辑。

\subsection{A5 关键细节五：verify() 五项数值校验体系}

\texttt{verify()} 把算法正确性归结为五个可独立检验的投影恒等式\textcolor{evidence}{\texttt{\nolinkurl{lib/multigrid.cpp:783-1069}}}：
(1) $0=(1-PP^\dagger)v_k$——零空间向量确实落在粗空间投影内（正交完备性）；
(2) $0=(1-P^\dagger P)\eta_c$——粗变量被 $P$ 无损嵌入；
(3) $0=(D_c-P^\dagger DP)\eta$——原生粗算子与模拟算子一致（Galerkin 性质）；
(4) $0=(\hat D_c-X^{-1}D_c)$——奇偶预条件形式正确（Deo/Doe 分别检验，\textcolor{evidence}{\texttt{\nolinkurl{multigrid.cpp:1001--1034}}}）；
(5) $|\mathrm{Im}\langle x|M^\dagger M|x\rangle|<\epsilon$——正规算子性质（Hermiticity 链路完好，\textcolor{evidence}{\texttt{\nolinkurl{multigrid.cpp:1054--1069}}}）。
这套内建自检是数值库少见的设计：任何转移算子/粗算子的回归缺陷都会在 setup 末尾立即暴露，而非等到求解阶段表现为收敛劣化。

其中第 (5) 项对 twisted-mass 前缀尤为关键：clover 项中的 $i\mu\gamma_5\tau^3$ 分量使算子的 $\gamma_5$-厄米链推广为逐层带 $\mu$ 缩放的广义形式，各层 \texttt{mu\_factor} 差值必须精确补偿（\textcolor{evidence}{\texttt{\nolinkurl{multigrid.cpp:371--376}}}、980--991）；一旦粗化破坏了该结构（例如缩放因子遗漏或符号错配），$M^{\dagger}M$ 不再严格厄米，以 $\langle\cdot,\cdot\rangle$ 为内积的 Krylov 求解器将丧失极小残差最优性，表现为 setup 后期或求解首段的收敛停滞。第 (5) 项以随机向量的虚部范数作哨兵，正是把这类``结构性回归''拦截在 setup 出口的最廉价手段。

\subsection{A6 任务如何正确执行（使用侧正确姿势）}

标准生命周期（tests/invert\_test.cpp:239--385 实证）：填 \texttt{mg\_param}（注意逐层数组都要赋值，\texttt{setMultigridParam} 提供模板）→ \texttt{newMultigridQuda(\&mg\_param)} 得预条件子 → \texttt{inv\_param.preconditioner = mg} 且 \texttt{inv\_type} 取 GCR，随后多次调用 \texttt{invertQUDA}（可穿插 \texttt{updateMultigridQuda} 更新规范场）→ \texttt{destroyMultigridQuda(mg)}。易错点：(a) 外层 solve\_type 必须为 DIRECT\_SOLVE（interface\_quda.cpp:2869--2870 强制）；(b) 聚合尺寸必须整除格点维度且粗维为偶数，否则 Transfer 构造期警告/减半重试（\textcolor{evidence}{\texttt{\nolinkurl{transfer.cpp:41--54}}}）；(c) \texttt{n\_vec} 不得超过聚合自由度（\textcolor{evidence}{\texttt{\nolinkurl{transfer.cpp:82--83}}}）。

\subsection{A7 实际执行情况（本次分析的实测记录）}

本次为静态只读分析，未编译运行 GPU 程序（环境无 CUDA 目标硬件需求声明，且技能约束为只读）。实测动作与产出：grep/read 定位 14 个核心文件的行级证据并逐一核实（引用前二次核对行号，如 V-cycle 段 1186--1213、$R$ 核 92--116）；wc 统计核心文件规模（multigrid.cpp 1710 行、coarse\_op.cuh 1465 行、multigrid.h 753 行等，合计 6532 行）；pdftotext 检索 docs/ 三份历史 PDF 命中 multigrid 条目 29 处。全程命令与退出码记录于会话日志 \texttt{\nolinkurl{.auto.2026-08-24-11-34-51.log}} 与 \texttt{\nolinkurl{.analy.2026-08-24-*.log}}。动态性能数字（如加速比）\textcolor{warn}{未验证}——本报告不做任何未经运行的性能断言。

\subsection{A8 实际输出情况}

本次任务产出物清单：（1）本报告 \texttt{\nolinkurl{docs/analy_multigrid_20260824.pdf}} 及同源 .tex；（2）会话日志两份（auto 与 analy 各一，不入库）；（3）用户输入存档 \texttt{\nolinkurl{.agent.2026-08-24-11-29-43.list}}。无二进制/数据产物（只读约束下唯一写操作即 docs/ 与日志，符合技能边界）。

\subsection{A9 结果分析（证据链综合解读）}

将收集到的 40 余处行级证据按算法主线串联，可得如下量化画像：setup 递归的复杂度集中在 \texttt{reset()}（126 行）与其调用的三个 create*（合计约 360 行），而求解热路径 \texttt{operator()} 仅 94 行——\textbf{结构上``一劳永逸的重 setup''换取``极轻的 per-step 开销''}，这与多重网格``setup 成本摊销到上千次 solve''的使用模式匹配。行数分布也印证了工程重心：粗算子构建族（coarse\_op*.cuh/.in.* 约 4000+ 行含实例化）远大于驱动逻辑本身，说明 Galerkin 投影的 GPU 实现才是复杂度主体。三份历史报告的独立结论（``近零模聚合实现粗化是本库最具辨识度的算法''，\textcolor{reference}{pure\_quda\_core\_20260815 报告 C1 节}）与本报告从源码独立推出的结论一致，构成旁证闭环。

\subsection{A10 综合分析：代码—对象映射与深层原因}

\begin{xltabular}{\textwidth}{llX}
\toprule
代码符号 & 对象概念 & 物理/数学含义 \\
\midrule
\endhead
$B$（\texttt{std::vector<ColorSpinorField>}） & 近零模向量集 $\{\psi_k\}_{k=1}^{N_{vec}}$ & $D\psi\approx0$ 的平滑解，编码算子红外结构（长程物理） \\
\texttt{geoBlockSize} & 聚合 aggregate & 相邻 $2^4$ 格点捆绑为一粗格点 \\
$V$ 场 & 延拓分量矩阵 & 块内正交基 $\phi_k$，\texttt{nColor}$\times$\texttt{Nvec} 打包 \\
$P$/\texttt{Transfer::P} & 延拓算子 & 式\eqref{eq:prolongate}，粗解插值回细格 \\
$R$/\texttt{Transfer::R} & 限制算子 & 式\eqref{eq:restrict}，$R=P^\dagger$ 块正交投影 \\
$Y_\mu$ & 粗规范链接 & 规范协变的聚合间耦合 $\sum\overline{V}UV$ \\
$X$ & 粗 clover & $P^\dagger DP$ 的接触项（对角块） \\
$\hat Y=X^{-1}Y$ & 预条件粗链接 & 使粗 dslash 免除显式求逆 \\
\texttt{presmooth/postsmooth} & 高频滤波器 & MR/GCR($+$additive Schwarz) $\nu_{pre},\nu_{post}$ 次 \\
\texttt{coarse\_solver} & 低频精确解 & 粗格 Krylov 至 \texttt{coarse\_solver\_tol} \\
\texttt{cycle\_type} & 环路拓扑 & V/F/W/递归，误差频率分离策略 \\
\texttt{mu\_factor} & twisted 质量缩放 & 逐层 $\mu\to f\cdot\mu$ 维持 twist 一致性 \\
\texttt{spin\_map} & 手征映射 & Wilson 4 旋量 $\to$ 2 粗旋量（左右手征块） \\
\bottomrule
\caption{代码符号↔物理对象映射表（数据来源：源码实测）}
\end{xltabular}

深层原因层面：为何 $R$ 取 $P^\dagger$ 而非几何 MG 的加权平均？——因为粗空间是\textbf{数值定义}的子空间，只有取厄米共轭才能保证 Galerkin 投影保持 $D$ 的对称化结构（$D^\dagger D$ 的正定性、$D$ 的 $\gamma_5$ 厄米性），这是 verify 第 3/5 项成立的前提；为何粗算子要显式构建而非每次 apply 时现算 $RDP$？——setup 一次构建后粗 apply 成本与普通 dslash 同阶，否则每次残差下行都要付出三次全场收缩。这两点是``代数''二字在性能与正确性上的双重兑现。

\subsection{A11 结尾总结}

\begin{conclbox}
QUDA MG 的整体算法思路 = \textbf{聚合代数多重网格 + 自适应近零空间 + Galerkin 投影粗化}；整体结构 = \textbf{每层一个 MG 组合体（Transfer + 粗 Dirac + 双平滑器 + 粗解器）的递归对象树}，setup 期自顶向下生长、solve 期自底向上递归。``自适应''由 generateNullVectors 的生成—投影—重建闭环承担，正确性由 verify() 五项投影恒等式守门。
\end{conclbox}

\subsection{A12 结尾评价}

优点：（1）结构清晰，层次职责单一，递归树天然支持任意层数（上限 5）；（2）数值纪律严格——混合精度链、固定点两遍块正交（\texttt{blockOrthoTwoPass}）、verify 内建校验；（3）可扩展性好——staggered KD、MMA、CA 求解器均为插件式增量。可改进处：（1）FCYCLE/WCYCLE 未在粗解器层显式实现（依赖推断行为）；（2）\texttt{MG} 构造/析构的手工 new/delete 链较长（\textcolor{evidence}{\texttt{\nolinkurl{multigrid.cpp:704--735}}}），异常安全性依赖调用纪律；（3）双精度 MG 需编译期 \texttt{GPU\_MULTIGRID\_DOUBLE} 显式开启（\textcolor{evidence}{\texttt{\nolinkurl{multigrid.h:10--13}}}）。

\subsection{A13 补充说明（假设与局限）}

本分析假设快照与上游 QUDA 1.1.0 功能等价（版本号与 NEWS 一致，但该快照处于 PyQCU 参考仓共享 git 下，\textcolor{warn}{未验证}上游 diff）。未覆盖：madwf\_transfer（Domain-wall MG 扩展，仅登记）、CPU 路径（\texttt{QUDA\_CPU\_FIELD\_LOCATION} 分支未深入）、QMP/MPI 通信细节。性能结论一律缺省（见 A7）。

\subsection{A14 参考源}

见第~\ref{sec:refs}~节参考源清单（编号连续引用）。

\subsection{A15 附录}

原始证据以日志形式存档：\texttt{\nolinkurl{.analy.2026-08-24-*.log}}（调查命令与证据清单）、\texttt{\nolinkurl{.auto.2026-08-24-11-34-51.log}}（auto 编排记录），头部摘录见第~8~节。

% ================= 6 关键代码片段 =================
\section{关键代码片段（直引原文件）}

以下片段全部以 \texttt{\textbackslash lstinputlisting} 直引仓库原文件，与源码逐字节一致。

\subsection{V-cycle 主循环核心段}
\lstinputlisting[firstline=1186,lastline=1211]{/root/PyQCU/refer/git-rep/quda/lib/multigrid.cpp}

第 1193 行限制、1196 行粗解、1200--1206 行延拓累加，即式\eqref{eq:restrict}与\eqref{eq:prolongate}的调度点。

\subsection{限制核 $R=P^\dagger$（块内归约投影）}
\lstinputlisting[firstline=92,lastline=116]{/root/PyQCU/refer/git-rep/quda/include/kernels/restrictor.cuh}

\texttt{cmac(conj(arg.v(...)), in(...), ...)} 即式\eqref{eq:restrict}的被积项；外层 \texttt{BlockReduce} 完成聚合内求和。

\subsection{延拓核 $P$（色空间旋转）}
\lstinputlisting[firstline=96,lastline=132]{/root/PyQCU/refer/git-rep/quda/include/kernels/prolongator.cuh}

\subsection{自适应 setup：生成—正交化主循环}
\lstinputlisting[firstline=1365,lastline=1392]{/root/PyQCU/refer/git-rep/quda/lib/multigrid.cpp}

\subsection{自适应 setup：重建闭环（递归 reset）}
\lstinputlisting[firstline=1429,lastline=1452]{/root/PyQCU/refer/git-rep/quda/lib/multigrid.cpp}

第 1438--1442 行把新零空间限制下传，1444--1445 行触发粗层重建——自适应性的落点。

\subsection{verify 第一项：$0=(1-PP^\dagger)v_k$}
\lstinputlisting[firstline=783,lastline=798]{/root/PyQCU/refer/git-rep/quda/lib/multigrid.cpp}

\subsection{setup 递归头部（reset）}
\lstinputlisting[firstline=72,lastline=100]{/root/PyQCU/refer/git-rep/quda/lib/multigrid.cpp}

\subsection{聚合查找表生成}
\lstinputlisting[firstline=214,lastline=237]{/root/PyQCU/refer/git-rep/quda/lib/transfer.cpp}

\subsection{外层 GCR-MG 包装（求解器工厂）}
\lstinputlisting[firstline=69,lastline=76]{/root/PyQCU/refer/git-rep/quda/lib/solver.cpp}

% ================= 7 图表详览 =================
\section{本次大任务图表详览}

本次大任务为纯静态只读分析，\textbf{未产生独立图像文件}；全部结构信息以表格（表 1--5）与文本依赖链（第 3 节 keybox）呈现，此为图表清单的闭合声明。报告内嵌的``准图表''及其来源如下表。

\begin{xltabular}{\textwidth}{llX}
\toprule
编号 & 图表载体 & 详尽介绍（来源/含义/解读/关联） \\
\midrule
\endhead
F1 & 表 1（目录职责表） & 来源：find 实测计数；含义：仓库六层物理布局；解读：lib+include 占 79\% 文件量，佐证``实现密集型''库定位；关联：第 2 节分层结构的物理基础 \\
F2 & 表 2（统计表） & 来源：find/wc/git 实测；含义：语言构成与 MG 深度常量；解读：CUDA/C++ 双轨；关联：A1 的层数上限论证 \\
F3 & 表 3（覆盖表） & 来源：四路径遍历登记；含义：饱和覆盖证明；解读：精读 5+浏览 20+仅索引余量，无死角；关联：方法论合规性 \\
F4 & 表 4（生命周期表） & 来源：\textcolor{evidence}{\texttt{\nolinkurl{multigrid.cpp:16--1224}}} 通读归纳；含义：setup/solve/reset 三阶段动作序列；解读：重 setup 轻 solve 的结构倾斜；关联：A9 量化画像 \\
F5 & 表 5（对象映射表） & 来源：A10 综合；含义：代码符号↔物理对象双语对照；解读：13 项映射无一推测（均有行号锚点）；关联：全文物理图像收束 \\
\bottomrule
\caption{本次大任务图表清单（穷尽闭合；无外部图像文件，数据来源：实测）}
\end{xltabular}

% ================= 8 日志汇编 =================
\section{本次大任务日志关键内容汇编}

本次大任务产生的全部日志如下表；头部片段直引于后（原始全文不入库）。

\begin{xltabular}{\textwidth}{llX}
\toprule
编号 & 日志文件 & 详尽介绍（用途/关键条目/解读/关联） \\
\midrule
\endhead
L1 & \texttt{\nolinkurl{.auto.2026-08-24-11-34-51.log}} & auto 编排日志：会话头（目标技能 analy、任务定义）、Step3--4 调查完成分隔块（覆盖表闭合声明与 12 条关键证据行号）；支撑 A7/A8 的实测记录主张 \\
L2 & \texttt{\nolinkurl{.analy.2026-08-24-*.log}} & analy 会话日志：全库调查命令与索引统计、证据清单逐条 文件:行号、tex 生成与编译结果、PDF 验证；是 A7 ``实测动作''的一手凭证 \\
L3 & \texttt{\nolinkurl{.agent.2026-08-24-11-29-43.list}} & 用户输入存档（2 条）：身份设定与本任务指令；界定分析范围与``着重算法思路与结构''的要求来源 \\
\bottomrule
\caption{本次大任务日志清单（穷尽闭合，原始片段见下方）}
\end{xltabular}

\lstinputlisting[language=,firstline=1,lastline=8]{/root/PyQCU/refer/git-rep/quda/.auto.2026-08-24-11-34-51.log}

% ================= 9 参考源清单 =================
\section{参考源清单}\label{sec:refs}

\begin{xltabular}{\textwidth}{llX}
\toprule
\textcolor{evidence}{参考源} & 类型 & 内容/作用 \\
\midrule
\endhead
\texttt{\nolinkurl{include/multigrid.h:81-258}} & 仓库 & MGParam 逐层元数据结构（A2 参数输入的结构化形态） \\
\texttt{\nolinkurl{include/multigrid.h:263-495}} & 仓库 & MG 类成员与接口（递归树骨架） \\
\texttt{\nolinkurl{include/multigrid.h:433-443}} & 仓库 & verify 五项恒等式文档注释 \\
\texttt{\nolinkurl{include/multigrid.h:585-687}} & 仓库 & CoarseOp/CoarseCoarseOp 声明（Galerkin 构建 API） \\
\texttt{\nolinkurl{include/quda_constants.h:46-50}} & 仓库 & QUDA\_MAX\_MG\_LEVEL=5 \\
\texttt{\nolinkurl{include/quda.h:636-830}} & 仓库 & QudaMultigridParam 全参数表 \\
\texttt{\nolinkurl{include/enum_quda.h:178-184}} & 仓库 & cycle 类型枚举 V/F/W/RECURSIVE \\
\texttt{\nolinkurl{include/enum_quda.h:476-482}} & 仓库 & transfer 类型枚举 AGGREGATE/KD 族 \\
\texttt{\nolinkurl{include/transfer.h:30-220}} & 仓库 & Transfer 类：V 场、双向查找表、P/R 接口 \\
\texttt{\nolinkurl{lib/multigrid.cpp:16-70}} & 仓库 & MG 构造：零空间来源四分支 \\
\texttt{\nolinkurl{lib/multigrid.cpp:72-197}} & 仓库 & reset():setup 递归主体 \\
\texttt{\nolinkurl{lib/multigrid.cpp:273-340}} & 仓库 & createSmoother:pre/post 平滑器装配 \\
\texttt{\nolinkurl{lib/multigrid.cpp:342-434}} & 仓库 & createCoarseDirac:粗算子三件套装配 \\
\texttt{\nolinkurl{lib/multigrid.cpp:564-702}} & 仓库 & createCoarseSolver 的 VCYCLE 与 RECURSIVE 分派及粗 deflation \\
\texttt{\nolinkurl{lib/multigrid.cpp:745-1129}} & 仓库 & verify() 实现 \\
\texttt{\nolinkurl{lib/multigrid.cpp:1131-1224}} & 仓库 & operator():V-cycle \\
\texttt{\nolinkurl{lib/multigrid.cpp:1275-1473}} & 仓库 & generateNullVectors:自适应引擎 \\
\texttt{\nolinkurl{lib/transfer.cpp:22-115}} & 仓库 & Transfer 构造：块校验/聚合约束/映射表 \\
\texttt{\nolinkurl{lib/transfer.cpp:152-164}} & 仓库 & BlockOrthogonalize 调用（V 生成） \\
\texttt{\nolinkurl{include/kernels/restrictor.cuh:92-116}} & 仓库 & R=$P^\dagger$ 投影核 \\
\texttt{\nolinkurl{include/kernels/prolongator.cuh:96-132}} & 仓库 & P 延拓核 \\
\texttt{\nolinkurl{include/kernels/coarse_op_kernel.cuh:1074-1114}} & 仓库 & VUV 块稀疏收缩与 $\gamma_\mu$ 合并 \\
\texttt{\nolinkurl{lib/dirac_coarse.cpp:461-479}} & 仓库 & createCoarseOp 聚合约束 \\
\texttt{\nolinkurl{lib/solver.cpp:69-76}} & 仓库 & GCR-MG 工厂包装 \\
\texttt{\nolinkurl{lib/interface_quda.cpp:2853-2939}} & 仓库 & multigrid\_solver 构造与 newMultigridQuda \\
\texttt{\nolinkurl{tests/utils/set_params.cpp:361-491}} & 仓库 & setMultigridParam 默认配置 \\
\texttt{\nolinkurl{tests/invert_test.cpp:239,385}} & 仓库 & MG 生命周期调用点 \\
\texttt{\nolinkurl{README.md:28-32}} & 仓库 & 官方定位陈述 \\
\texttt{\nolinkurl{NEWS:29-35}}, \texttt{\nolinkurl{NEWS:207}} & 仓库 & 演进史（初版实现→tensor core/staggered 增量） \\
\textcolor{reference}{\texttt{\nolinkurl{docs/pure_quda_core_20260815.pdf}}} & 参考 & 历史剖析报告：MG 深挖章（C1），结论与本研究互证 \\
\textcolor{reference}{\texttt{\nolinkurl{docs/pure_quda_20260817.pdf}}} & 参考 & 历史剖析报告：自适应 MG 定位为独特竞争力之首 \\
\textcolor{reference}{\texttt{\nolinkurl{docs/analy_quda_20260817.pdf}}} & 参考 & 历史全库分析：multigrid 文件族索引背景 \\
\bottomrule
\caption{参考源清单（类型``仓库''为直接证据，``参考''为知识库旁证；数据来源：read/grep/pdftotext 实测）}
\end{xltabular}

% ================= 10 结论与局限 =================
\section{结论与局限}

\begin{conclbox}
三视角收束。\textbf{结构}：六层架构（入口/声明/驱动/转移/粗算子/工厂）之上，\texttt{MG} 类以组合模式构成递归树，每层自持全部构件。\textbf{关系}：构造期依赖注入（Transfer/$Y$/$X$/smoothers 均在 reset 中挂接）+ 运行期递归调用（V-cycle 三次跨界：$R$ 下行、粗解、$P$ 上行）。\textbf{思路}：以数值近零空间定义粗空间的聚合 AMG，用 Galerkin 投影保结构，用生成—投影—重建闭环实现自适应，用内建 verify 保正确。
\end{conclbox}

\begin{refbox}
知识库旁证：本库三份既有分析报告一致认定自适应多重网格为 QUDA 独特竞争力之首（\textcolor{reference}{pure\_quda\_core\_20260815 C1 节}``近零模聚合实现粗化，是本库最具辨识度的算法''），与本研究从源码独立导出的结论吻合。
\end{refbox}

\begin{warnbox}
局限与未验证项：（1）快照与上游官方 1.1.0 的差异 \textcolor{warn}{未验证}（无上游比对基线）；（2）FCYCLE/WCYCLE 的复合实现路径为\textcolor{warn}{推断}级别（enum 存在但粗解器层未见独立分支）；（3）动态性能（加速比、setup 耗时占比）\textcolor{warn}{未验证}——静态分析不作性能断言；（4）madwf（Domain-wall MG）、MPI halo 细节仅登记未深挖；（5）docs/ 三份 PDF 为二手旁证，其行号引用未逐条复核（仅主题结论互证）。
\end{warnbox}


% ================= 附录：QUDA ↔ PyQCU 映射 =================
\section{\texorpdfstring{附录：QUDA $\leftrightarrow$ PyQCU 概念映射}{附录：QUDA 与 PyQCU 概念映射}（2026-08-24 增补）}

以下映射基于 PyQCU dev84 攻坚期的对照实践（详见 \texttt{logs/fix-report-2026-08-24.md}
与 \texttt{examples/qcu/dev84/dev84\_report.md}），供跨库阅读时定位对应实现。

\begin{center}
\begin{tabular}{p{0.44\textwidth}p{0.44\textwidth}}
\hline
\textbf{QUDA 1.1.0} & \textbf{PyQCU} \\
\hline
\texttt{generateNullVectors} (\textcolor{evidence}{\texttt{\nolinkurl{multigrid.cpp:1275}}}) & \texttt{tools.give\_null\_vecs} / \texttt{give\_null\_vecs\_mt}（逆迭代 $v \leftarrow v - A^{-1}Av$，相对容差） \\
\texttt{transfer->R} / \texttt{transfer->P} & \texttt{tools.restrict} / \texttt{tools.prolong}；CUDA 路径 \texttt{applyMultigridRestrictQcu} / \texttt{applyMultigridProLongQcu} \\
Galerkin 粗算子构建（coarse\_op.cuh VUV 族） & 33-tensor stencil：\texttt{build\_stencil} / \texttt{build\_stencil\_local}（BatchedLocalSchur 窗口化）+ \texttt{apply\_stencil} \\
\texttt{MG::operator()} V-cycle & \texttt{solver.multigrid.cycle / solve}；多线程多卡 \texttt{MultiGpuMultigrid} \\
smoother setup（pre/post，return\_residual） & 层 0 C++ BiStabCG 平滑（\texttt{applyCloverBistabCgDslashQcu}）；层 2+ Python \\
粗解器（coarse solver） & 粗层 BiStabCG（\texttt{\_coarse\_dslash\_cuda} + \texttt{bistabcg}）/ Python einsum \\
柔性外迭代（FGMRES/GCR 预条件 MG） & \texttt{solver.fgmres(precond=V-cycle)}；dev84 门控 GCR 路径 \\
verify() 五项投影恒等式 & \texttt{pyqcu.testing.verify\_nullvecs}（近零性/正交性/Galerkin 一致性抽样） \\
nullvec 缓存与测试向量 & \texttt{data/*.h5} 缓存（lonv/hnn/hdg/sit 单句柄一次写全） \\
\hline
\end{tabular}
\end{center}

\noindent\textcolor{reference}{差异提示：PyQCU 的奇偶子格沿 $t$ 维紧缩（\texttt{lat\_shape} 返回 $T/2$），
与 QUDA 的逐维聚合块布局不同；跨库移植转移算子时须先对齐坐标约定。}

\end{document}